\chapter{Woodwinds and Brass: Speak with Colour, Answer with Accents}
\label{ch:winds-brass}

\classquote{Woodwinds are not ``lighter strings'', and brass is not ``louder strings''. One handles colour and line, the other handles declaration and energy; use the wrong role and even a good patch sounds pasted on.}{}{}

\section{Two families, two jobs}
\label{sec:winds-jobs}

The \term{woodwind section}{woodwind-section} includes flute, oboe, clarinet and bassoon; the \term{brass section}{brass-section} includes trumpet, horn, trombone and tuba. Their most common division of labour is:

\begin{itemize}
\item \textbf{Woodwinds}: recolour a melody, add lines to the harmony, answer where the vocal leaves space; their strength is agility, detail and strongly differentiated timbres.
\item \textbf{Brass}: ``make a statement'' at section changes, chorus entries and endings; their strength is projection, weight and ceremony.
\item \textbf{Neither should live everywhere}: woodwinds through an entire song become noise, brass through an entire song becomes shouting; both work better as entrances and exits than as a permanent bed.
\end{itemize}

Before writing either family, ask three questions: does this passage need colour or energy? On which beat should it enter? Does it fight the vocal for the same register? Answer those before choosing instruments.

\section{Woodwinds: four colours, four lines}
\label{sec:woodwinds}

Each woodwind has a very different timbre; they cannot be stacked like a string section:

\begin{itemize}
\item \textbf{Flute}: high register, light line, suited to fast runs, high counter-melodies and bright colour; do not park it in the low-middle register.
\item \textbf{Oboe}: focused, nasal, strongly projecting; suited to solo-like long phrases and a ``weeping'' colour, not dense accompaniment.
\item \textbf{Clarinet}: flexible in the middle, rounded in colour, suited to connecting inner voices and duetting with the flute; leave space in the low register.
\item \textbf{Bassoon}: the low woodwind, holding the section's floor; suited to bass lines, secondary melodies and reedy colour, not high-speed runs.
\end{itemize}

The key to writing woodwinds is that \textbf{every line has an identity}: flute sings the high counter-melody, clarinet fills the middle, bassoon holds the bass, oboe appears only in the key phrase. If all four play the lead, the result is not ``thick'' but ``four versions of the lead fighting''.

\section{Brass: declaration, accent, weight}
\label{sec:brass}

Brass works more like punctuation than background harmony:

\begin{itemize}
\item \textbf{Trumpet}: the top of the brass choir, responsible for declarations, openings and the first beat of the chorus; bright, and a long continuous melody will overpower the vocal.
\item \textbf{Horn}: round and full in the middle, suited to inner voices and blending with cello/viola; it is the easiest brass instrument to fold into the texture.
\item \textbf{Trombone}: tenor/bass, suited to roots, long notes and cadences; slides and short phrases add weight, but do not write it as a second bass.
\item \textbf{Tuba}: the lowest voice, holding the root; not a melody instrument, and not to be forced into quiet passages.
\end{itemize}

A \term{brass stab}{brass-stab} is the most effective brass device: a one-bar or two-beat short chord on the first beat of the chorus, a section change or the final phrase. It does not need continuous playing; it only needs to ``stamp'' the key moment. Too many stabs become noise every bar; too few and the brass disappears.

\section{Register and projection}
\label{sec:winds-ranges}

The following table gives conservative working ranges:

\begin{longtable}{p{0.18\textwidth}p{0.20\textwidth}p{0.42\textwidth}}
\toprule
Instrument & Working range & Orchestral use \\
\midrule
\endfirsthead
\toprule
Instrument & Working range & Orchestral use \\
\midrule
\endhead
Flute & C4--C7 & High counter-melody, fast runs, bright colour \\
Oboe & Bb3--F6 & Solo long phrases, projecting colour; rarely dense chords \\
Clarinet & D3--F6 & Inner connection, duet with flute \\
Bassoon & Bb1--Bb4 & Woodwind bass, secondary melody; no fast runs \\
Trumpet & E3--Bb5 & Declaration, accent, top of the brass choir \\
Horn & F3--C6 & Inner voice, blend with strings, round colour \\
Trombone & E2--Bb4 & Root, long notes, cadence; not a second bass \\
Tuba & E1--Bb3 & Lowest root; not a melody \\
\bottomrule
\end{longtable}

\term{Projection}{projection} is not volume: oboe and trumpet ``cut through'' more than flute or horn at the same fader setting. Solve it first with register and rhythm in the arrangement; do not wait for the mix to force it down with a fader. Brass as a whole covers a vocal more easily than woodwinds, especially when the trumpet sits in the vocal's register.

\section{Combination: family, mixed, and avoiding the vocal}
\label{sec:winds-combination}

\term{Family grouping}{family-grouping} pairs woodwinds with woodwinds and brass with brass; \term{mixed grouping}{mixed-grouping} uses both. Common approaches:

\begin{enumerate}
\item \textbf{Woodwind duet}: flute and clarinet, one high and one middle, playing two independent lines; separate their registers, do not double in unison.
\item \textbf{Woodwind quartet}: flute, oboe, clarinet and bassoon each with a line; use it only briefly in the chorus or bridge and leave the verse empty.
\item \textbf{Brass choir}: trumpet, horn, trombone and tuba arranged high--middle--middle-low--low; do not let all four play the lead.
\item \textbf{Woodwind answer}: after the vocal phrase, two bars of woodwind reply; the answer does not repeat the vocal, it catches the emotion.
\item \textbf{Brass stab}: a short chord on the chorus downbeat or a section change; exit immediately and give the space back to the vocal.
\end{enumerate}

Use \term{doubling}{doubling} carefully: handing the lead unchanged to trumpet or flute covers the vocal's register and colour. If doubling is needed, choose an octave below, a third/sixth, or double only the final chorus phrase; never keep woodwinds/brass on the same notes and rhythm as the vocal for the whole song.

\begin{figure}[htbp]
  \centering
  \begin{tikzpicture}[font=\small]
    \draw[->] (0,0) -- (9.2,0) node[right] {bars};
    \foreach \x in {1,...,8} {\draw[gray!30] (\x,0.1) -- (\x,3.3); \node[font=\scriptsize] at (\x,0.05) {\x};}
    \fill[yebanblue!15] (0.2,2.35) rectangle (6.2,2.9);
    \node[font=\scriptsize] at (3.2,2.62) {vocal phrase};
    \fill[yebanblue!55] (2.2,1.45) rectangle (3.2,1.9);
    \fill[yebanblue!55] (4.2,1.45) rectangle (5.2,1.9);
    \node[font=\scriptsize] at (2.7,2.12) {WW answer};
    \node[font=\scriptsize] at (4.7,2.12) {WW answer};
    \fill[orange!55] (6.2,0.55) rectangle (7.2,1.0);
    \fill[orange!55] (7.2,0.55) rectangle (8.2,1.0);
    \node[font=\scriptsize] at (6.7,0.77) {brass};
    \node[font=\scriptsize] at (7.7,0.77) {brass};
    \node[font=\scriptsize, anchor=west] at (8.4,1.67) {woodwinds};
    \node[font=\scriptsize, anchor=west] at (8.4,0.77) {brass};
  \end{tikzpicture}
  \caption{Woodwind answers and brass stabs across eight bars: the vocal occupies the first six, woodwinds answer only at phrase ends, and brass stamps the chorus turn.}
  \label{fig:winds-brass-entry}
\end{figure}

\section{Realising it in Yeban}
\label{sec:winds-yeban}

\begin{enumerate}
\item \textbf{One family member per track}: create one MIDI track for flute, oboe, clarinet and bassoon, and one each for trumpet, horn, trombone and tuba. Do not pack all woodwinds into one track where balance becomes impossible.
\item \textbf{Load SFZ instruments}: use File $\rightarrow$ Load SFZ instrument library... to load the matching patch on each track; one instrument per track so the timbres do not fight.
\item \textbf{Write short phrases in the piano roll}: press P / F7, write the two-bar woodwind answer first, then the one-bar brass stab; do not turn the whole section into continuous chords.
\item \textbf{Use the velocity lane for accents}: give the first note of a brass stab full velocity and let the rest fall away quickly; taper the woodwind answer at its end so it does not fight the vocal's last beat.
\item \textbf{Use the automation lane for entrances}: expand the automation lane (A) in the arrangement area, keep woodwinds very low in the verse and lift them for the chorus answer; open brass only for the stab bar.
\item \textbf{Finish with pan and faders}: in the mixer (X / F8), place woodwinds and brass at different pan positions; solve balance in the arrangement first, then use faders.
\item \textbf{Check with Solo/Mute}: solo the vocal and confirm woodwinds/brass do not cover it with the same notes and rhythm; then solo the brass and confirm the stab does not bury the drum turn.
\end{enumerate}

\screenshot{arrangement-winds-brass}{Woodwind and brass tracks, piano roll and automation lane on one screen: short woodwind answers and short brass stabs with visible entrances}

\section{Common mistakes and fixes}
\label{sec:winds-mistakes}

\begin{longtable}{p{0.34\textwidth}p{0.58\textwidth}}
\toprule
Common mistake & Fix \\
\midrule
\endfirsthead
\toprule
Common mistake & Fix \\
\midrule
\endhead
Woodwinds run through the whole song & Answer only where the vocal leaves space; keep the verse empty and enter with short chorus phrases. \\
Brass sounds every bar & Reduce it to a short stab at a section change or the chorus downbeat, then exit. \\
All four woodwinds play the lead & Assign a high counter-line, inner voice, bass and solo colour; every line needs an identity. \\
Trumpet sits in the vocal's register & Move the trumpet down an octave or use it only for short stabs; leave the vocal register open. \\
All four brass instruments play chords & Divide high--middle--middle-low--low; do not treat all four as melody. \\
Bassoon is written fast runs & Give fast lines to flute/clarinet; keep bassoon on bass and long phrases. \\
Woodwinds/brass double the vocal in unison & Use an octave below, a third/sixth, or double only the final chorus phrase. \\
\bottomrule
\end{longtable}

\section{Exercise}
\label{sec:winds-exercise}

\begin{exercise}
Write eight bars of ``vocal + woodwind answer + brass stab'':
\begin{enumerate}
\item Write the vocal's four-bar phrase first, leaving two beats empty at the end; woodwinds answer only in that space with two bars.
\item Use flute and clarinet for the woodwind answer, with separated registers; do not copy the vocal phrase.
\item In bar 7 write a brass stab: one note each for trumpet, horn, trombone and tuba, arranged high--middle--middle-low--low; let it sound for only a beat and a half.
\item Use the velocity lane so the stab enters cleanly and exits cleanly; then use the automation lane to make the woodwinds almost inaudible in the verse.
\end{enumerate}
When finished, solo the vocal: do woodwinds and brass appear only where they should? If the music is clearer without them, the arrangement has already done their job.
\end{exercise}

\section{Chapter takeaways}
\label{sec:winds-takeaways}

\begin{itemize}
\item Woodwinds handle colour and line; brass handles declaration and energy; neither should live everywhere.
\item Woodwind timbres differ greatly; write separate lines instead of stacking the family blindly.
\item Brass is most effective in short stabs, not continuous chords.
\item Avoid the vocal first with register and rhythm; faders are the last patch.
\item Family grouping seeks blend, mixed grouping seeks contrast; doubling needs a new register or interval.
\item In Yeban: one instrument per track, separate SFZ patches, velocity/automation for entrances, Solo/Mute for acceptance.
\end{itemize}

\begin{helpmapping}
\begin{longtable}{p{0.45\textwidth}p{0.45\textwidth}}
\toprule
Section & Help anchor \\
\midrule
Two families, two jobs & \texttt{arrangement-winds-jobs} \\
Woodwinds: four colours, four lines & \texttt{arrangement-woodwinds} \\
Brass: declaration, accent, weight & \texttt{arrangement-brass} \\
Register and projection & \texttt{arrangement-winds-ranges} \\
Combination & \texttt{arrangement-winds-combination} \\
Realising it in Yeban & \texttt{arrangement-winds-yeban} \\
Common mistakes and fixes & \texttt{arrangement-winds-mistakes} \\
Exercise & \texttt{arrangement-winds-exercise} \\
\bottomrule
\end{longtable}
\end{helpmapping}
